\subsection{Relation to the literature}
Neural Bellman Operators build on numerical dynamic programming and function approximation. Projection and related methods exploit economic structure to approximate a Bellman equation; see \citet{judd1998}. Smooth neural networks can approximate derivatives under representation hypotheses \citep{hornik1990}. These results distinguish representational capacity from the accuracy achieved by a finite architecture and a particular optimizer.

The differential critic is related to the neural PDE approach in \citet{sirignano2018}. Their DGM formulation includes differential and boundary conditions and treats HJB equations. The deep BSDE method of \citet{han2018} also includes an HJB application. Our object combines a specified policy continuation, a feasible actor, and an economic error account, and asks when the continuation representation is useful across decision queries. An analytical maximizing action is useful when available; it is not a universal distinction between neural PDE methods. Recursive dependence on the value is a nonlinear zeroth-order term, whereas full nonlinearity of a PDE concerns its highest derivatives.

Recursive utility separates intertemporal substitution from risk aversion \citep{epsteinzin1989,duffieepstein1992}. Utility-process selection and transversality matter in its infinite-horizon formulation \citep{herdegen2021}. Temporal selves instead require an equilibrium formulation \citep{bjork2016}. The stochastic approximation literature provides convergence results under explicit step-size, stability, and limiting-dynamics conditions \citep{borkar2024}. We retain those distinctions when describing a finite neural computation.

Section~\ref{sec:model} defines the economic object and Section~\ref{sec:r10continuous} specializes it to the capital economy. Section~\ref{sec:r15mechanism} connects continuation error with economic gain. The computational design in Section~\ref{sec:r16design} distinguishes continuation reuse, conditional payoff replication, and mechanism assessment. Section~\ref{sec:economic-scope} states the other economic applications. General verification results and their full proofs remain in the technical supplement, beginning with Section~\ref{sec:theory}; the application companion is a current part of the paper.
